\section{Why Public Intervention? Private Evaluation and Social Welfare}
\label{sec:short_run}

Section~\ref{sec:hub_efficiency} characterized welfare under a benevolent planner for a stand-alone hub. For the paper's policy interpretation, however, it is also important to explain why a place-based public intervention may be needed in the first place. In practice, public hubs are often advocated or opposed using highly visible short-run indicators such as participation counts, meetings, first-round matches, or occupancy rates. These criteria do not necessarily coincide with social welfare.

This section therefore introduces a minimal contrast between a \emph{private evaluation rule} and a \emph{social evaluation rule}. The point is not to build a full decentralized equilibrium model. Rather, it is to show as parsimoniously as possible that local private stakeholders may evaluate a hub through visible short-run activity, whereas a social planner also values dynamic novelty and cross-regional spillovers. The two criteria need not coincide.

\subsection{A Private Evaluation Rule Based on Visible Short-Run Activity}

Suppose that a private hub operator, local coalition, or other decentralized stakeholder evaluates the hub primarily through visible short-run activity in region $A$. Under the present two-period setup, the natural reduced-form counterpart is the hub's period-1 net gain:
\begin{equation}
\Pi_H
:=
\bigl[B_1(1)-F\bigr]-B_1(0).
\label{eq:PiH}
\end{equation}
Using \eqref{eq:B1}, this becomes
\begin{equation}
\Pi_H
=
v_L P_A\bigl(m_A^1-m_A^0\bigr)-F.
\label{eq:DeltaSR}
\end{equation}

Expression \eqref{eq:DeltaSR} is the model's reduced-form representation of a private short-run evaluation rule. It depends only on the hub's immediate contribution to local matching in region $A$, net of the fixed cost. It ignores both dynamic redundancy in repeated local interaction and any social value created by future cross-regional linkages.

A hub is therefore \emph{privately supported under short-run evaluation} if and only if
\[
\Pi_H>0.
\]

\begin{proposition}
\label{prop:short_run_success}
Under the short-run private evaluation rule \eqref{eq:PiH}, a public hub is privately supported if and only if
\begin{equation}
F<v_L P_A\bigl(m_A^1-m_A^0\bigr).
\label{eq:short_run_condition}
\end{equation}
Equivalently, treating $N_A$ as continuous, private support obtains if and only if
\begin{equation}
N_A\ge N_{SR}^*
:=
\frac{1}{2}
\left(
1+
\sqrt{
1+\frac{8F}{v_L\bigl(m_A^1-m_A^0\bigr)}
}
\right),
\label{eq:NSRstar}
\end{equation}
with the usual ceiling operator if $N_A$ is restricted to integers.
\end{proposition}

Proposition \ref{prop:short_run_success} depends only on the local period-1 coordination gain. In that sense, it is deliberately myopic. This is useful because many real-world arguments for or against hubs are framed in exactly these terms.

\subsection{A Social Evaluation Rule}

The planner's criterion is broader. As shown in Sections~\ref{sec:hub_efficiency} and~\ref{sec:refresh}, social welfare depends not only on immediate local coordination, but also on dynamic redundancy in local interactions and on cross-regional renewal.

For notational convenience, define
\begin{equation}
S_A
:=
v_L P_A\bigl(m_A^1-m_A^0\bigr)>0,
\label{eq:SA}
\end{equation}
which is the hub's visible short-run local coordination gain,
\begin{equation}
D_A
:=
\beta v_L P_A\bigl(m_A^1-m_A^0\bigr)\Bigl[1-\delta\bigl(m_A^1+m_A^0\bigr)\Bigr],
\label{eq:DA}
\end{equation}
which is the dynamic correction generated by second-period local interaction, and
\begin{equation}
S_{AB}(b)
:=
\beta\eta v_X Q_{AB}\,b-\frac{\kappa}{2}b^2,
\label{eq:SAB}
\end{equation}
which is the incremental cross-regional renewal surplus generated by pipeline policy $b$.

Using \eqref{eq:DeltaH_repeat} and \eqref{eq:DeltaHb_repeat}, the hub-alone and hub-plus-pipeline welfare gains can be written as
\begin{equation}
\Delta_H
=
\Pi_H + D_A,
\label{eq:dynamic_wedge}
\end{equation}
and
\begin{equation}
\Delta_{H+b}(b)
=
\Pi_H + D_A + S_{AB}(b).
\label{eq:bundle_wedge}
\end{equation}

Equation \eqref{eq:dynamic_wedge} shows that the social value of a hub alone equals the private short-run criterion plus a dynamic correction term. Equation \eqref{eq:bundle_wedge} adds a third component: cross-regional renewal. The normative justification for public intervention is exactly that decentralized short-run evaluation omits these additional social terms.

Moreover, under the private short-run rule \eqref{eq:PiH}, pipeline policy has no direct private value because it affects neither period-1 local activity nor the private criterion itself. Hence a purely short-run private evaluator would set
\[
b^{P}=0,
\]
or at least would have no reason to choose a positive $b$. By contrast, the planner values $b$ through \eqref{eq:SAB} and chooses the positive optimum characterized in Section~\ref{sec:refresh}.

Let
\begin{equation}
R_{AB}(N_A,N_B)
:=
\frac{\bigl(\beta\eta v_X Q_{AB}\bigr)^2}{2\kappa}
=
\frac{\beta^2\eta^2 v_X^2 N_A^2N_B^2}{2\kappa}
>0
\label{eq:RAB_sec5}
\end{equation}
denote the maximized cross-regional renewal surplus from Section~\ref{sec:refresh}. Then
\begin{equation}
\Delta_{H+b}^*
=
\Pi_H + D_A + R_{AB}(N_A,N_B).
\label{eq:bundle_wedge_star}
\end{equation}

\subsection{When Do Private and Social Evaluations Diverge?}

The following proposition collects the two most important cases.

\begin{proposition}
\label{prop:dynamic_inefficiency}
Let
\[
F_H^*(N_A)=S_A+D_A
\qquad\text{and}\qquad
F_{H+b}^*(N_A,N_B)=S_A+D_A+R_{AB}(N_A,N_B).
\]

\begin{enumerate}[leftmargin=2em]
    \item \textbf{Private over-support of a hub alone.}
    If
    \begin{equation}
    \delta\bigl(m_A^1+m_A^0\bigr)>1,
    \label{eq:dynamic_gap}
    \end{equation}
    then $D_A<0$, so
    \[
    \Delta_H<\Pi_H.
    \]
    Moreover, there exists a nonempty interval of fixed costs for which the hub is privately supported under short-run evaluation but socially undesirable as a stand-alone policy:
    \begin{equation}
    \max\!\left\{0,\;F_H^*(N_A)\right\}
    <
    F
    <
    S_A.
    \label{eq:dynamic_inefficiency_interval}
    \end{equation}

    \item \textbf{Private under-support of a hub-plus-pipeline bundle.}
    If
    \begin{equation}
    D_A+R_{AB}(N_A,N_B)>0,
    \label{eq:bundle_gap_condition}
    \end{equation}
    equivalently if
    \[
    F_{H+b}^*(N_A,N_B)>S_A,
    \]
    then there exists a nonempty interval of fixed costs for which the hub is not privately supported under short-run evaluation, yet a hub combined with optimally chosen pipeline policy is socially desirable:
    \begin{equation}
    S_A
    \le
    F
    <
    F_{H+b}^*(N_A,N_B).
    \label{eq:under_support_interval}
    \end{equation}
\end{enumerate}
\end{proposition}

Proposition \ref{prop:dynamic_inefficiency} gives a compact rationale for place-based public intervention.

Part (1) shows that a hub may be \emph{privately attractive} because it generates visible local activity, yet \emph{socially undesirable} because repeated local interaction destroys too much future novelty. In this case, a purely decentralized short-run criterion overstates the hub's value.

Part (2) shows the opposite possibility. Once dynamic and cross-regional gains are taken into account, a hub combined with pipeline policy may be \emph{socially desirable} even though it does not pass the private short-run test. In this case, decentralized evaluation understates the value of the broader regional policy bundle.

Together, these two cases imply that the relevant policy question is not simply whether local actors can justify a hub using immediate activity metrics. The deeper question is whether the hub, viewed as part of a place-based coordination architecture, creates enough long-run and cross-regional value to justify public intervention.

\begin{figure}[htbp]
\centering
\begin{tikzpicture}
\begin{axis}[
    width=0.82\textwidth,
    height=0.52\textwidth,
    xmin=0, xmax=2.5,
    ymin=0, ymax=12.5,
    axis lines=left,
    xlabel={Local redundancy parameter $\delta$},
    ylabel={Hub fixed cost $F$},
    xtick={0,0.5,1.0,1.4286,2.0,2.5},
    xticklabels={$0$,$0.5$,$1.0$,$\delta^{\dagger}$,$2.0$,$2.5$},
    ytick={0,2,4,6,8,10,12},
    clip=false,
    legend style={draw=none, fill=none, font=\small, at={(0.02,0.98)}, anchor=north west},
    every axis plot/.append style={thick},
]

\def\PA{15}
\def\vL{1}
\def\mzero{0.15}
\def\mone{0.55}
\def\betaPar{0.9}

\pgfmathsetmacro{\DeltaM}{\mone-\mzero}
\pgfmathsetmacro{\SigmaM}{\mone+\mzero}
\pgfmathsetmacro{\SAplot}{\vL*\PA*\DeltaM}
\pgfmathsetmacro{\dcrit}{1/(\SigmaM)}

\addplot[name path=FH, domain=0:2.5, samples=200, black]
    {\vL*\PA*\DeltaM*((1+\betaPar)-\betaPar*x*\SigmaM)};
\addlegendentry{$F_H^*(\delta)$}

\addplot[name path=FSR, domain=0:2.5, samples=2, dashed, black]
    {\SAplot};
\addlegendentry{$S_A$ (short-run private threshold)}

\addplot[name path=zero, domain=0:2.5, samples=2, draw=none] {0};
\addplot[gray!18, draw=none] fill between[of=FH and zero];
\addplot[name path=FHright, domain=\dcrit:2.5, samples=200, draw=none]
    {\vL*\PA*\DeltaM*((1+\betaPar)-\betaPar*x*\SigmaM)};
\addplot[name path=FSRright, domain=\dcrit:2.5, samples=2, draw=none]
    {\SAplot};
\addplot[gray!45, draw=none] fill between[of=FSRright and FHright];
\addplot[densely dotted, black] coordinates {(\dcrit,0) (\dcrit,12.5)};

\node[font=\small, align=center] at (axis cs:0.60,2.7)
    {Hub privately supported\\and socially desirable};

\node[font=\small, align=center] at (axis cs:1.95,5.0)
    {Hub privately supported\\but socially undesirable};

\node[font=\small, align=center] at (axis cs:1.85,9.2)
    {No private support\\for hub};

\end{axis}
\end{tikzpicture}
\caption{Illustrative policy regions in $(\delta,F)$ space. The solid line shows the social efficiency threshold for a hub alone, $F_H^*(\delta)$, while the dashed line shows the short-run private support threshold, $S_A=v_LP_A(m_A^1-m_A^0)$. For $\delta>\delta^{\dagger}=1/(m_A^1+m_A^0)$, there exists a region in which the hub is privately supported under short-run evaluation but socially undesirable once dynamic redundancy is taken into account. Parameter values are illustrative: $N_A=6$, $v_L=1$, $m_A^0=0.15$, $m_A^1=0.55$, and $\beta=0.9$.}
\label{fig:delta_F_regions}
\end{figure}

\subsection{Interpretation}

The section delivers three implications.

First, a short-run private rule naturally ignores dynamic redundancy. This creates the possibility that local stakeholders favor a hub because it produces visible activity today, even though it lowers social welfare once repeated local matching is taken into account.

Second, a short-run private rule also ignores cross-regional spillovers. This creates the possibility that local stakeholders reject a hub because its visible short-run returns are too small, even though a hub combined with outside-linkage policy would be socially beneficial.

Third, the role of government in the present framework is therefore not merely to finance an otherwise identical private project. It is to evaluate the hub on a wider margin than decentralized actors do: not only by immediate local coordination, but also by dynamic renewal and by the creation of cross-regional pipelines. This is the paper's minimal theoretical rationale for public intervention.

Section~\ref{sec:refresh} builds directly on this logic by characterizing the optimal cross-regional pipeline policy and the conditions under which the efficient policy object is not a hub alone, but a hub combined with external linkage formation.
